% =====================================================================
\chapter{Forma polar, fórmula de Euler, De Moivre y rotaciones}
\label{tema:6}
% =====================================================================

La forma binómica es cómoda para sumar; la forma polar, para multiplicar. Este tema demuestra que al multiplicar dos complejos se multiplican los módulos y se suman los argumentos, introduce la exponencial imaginaria $e^{i\theta}$ y la fórmula de Euler, deduce el teorema de De Moivre y la fórmula de las raíces $n$-ésimas, y termina interpretando la multiplicación por $e^{i\alpha}$ como el giro de ángulo $\alpha$ del plano, que es la idea sobre la que se construye RoPE. Se siguen la sección~6 de las notas de 18.03 del MIT~\cite{mit1803}, las §1.2 a §1.5 de las notas de Cambridge~\cite{cambvm}, la §10.5 de~\cite{osat} y la lectura geométrica del capítulo~1 de Needham~\cite{needham}.

\section{Argumento y forma polar}

\begin{definicion}[Argumento y forma polar]\label{def:polar}
Sea $z\ne0$ y $r=\abs{z}$. Un \emph{argumento} de $z$ es cualquier $\theta\in\R$ tal que
\[
z=r(\cos\theta+i\sin\theta).
\]
Esta expresión es la \emph{forma polar} de $z$. El \emph{argumento principal} $\Arg z$ es el único argumento en $(-\pi,\pi]$. El número $0$ no tiene argumento.
\end{definicion}

Todo $z\ne0$ tiene argumento: $z/r$ tiene módulo~1, luego es un punto $P(\theta)$ de la circunferencia unidad. Por el \cref{teo:simetrias}(ii), dos argumentos de $z$ difieren en un múltiplo entero de $2\pi$, y exactamente uno está en $(-\pi,\pi]$.

\begin{proposicion}[Cálculo del argumento principal]\label{prop:arg}
Sea $z=a+bi\ne0$. Si $a>0$, $\Arg z=\arctan(b/a)$. Si $a<0$, $\Arg z=\arctan(b/a)+\pi$ cuando $b\ge0$ y $\arctan(b/a)-\pi$ cuando $b<0$. Si $a=0$, $\Arg z=\pi/2$ cuando $b>0$ y $-\pi/2$ cuando $b<0$.
\end{proposicion}

\begin{proof}
Si $\theta$ es un argumento, $\cos\theta=a/r$ y $\sin\theta=b/r$, luego $\tan\theta=b/a$ cuando $a\ne0$. Si $a>0$, $\theta$ está en el cuadrante I o IV y $\arctan(b/a)\in(-\frac{\pi}{2},\frac{\pi}{2})$ es el argumento correcto. Si $a<0$, $\arctan(b/a)$ apunta al cuadrante opuesto y hay que sumar o restar $\pi$ (\cref{teo:simetrias}(v)), eligiendo el signo que deja el resultado en $(-\pi,\pi]$. El caso $a=0$ es directo.
\end{proof}

\begin{ejemplo}\label{ej:polar}
Para $z=-1+\sqrt3\,i$: $r=\sqrt{1+3}=2$ y, como $a<0\le b$, $\Arg z=\arctan(-\sqrt3)+\pi=-\frac{\pi}{3}+\pi=\frac{2\pi}{3}$. Luego $z=2\bigl(\cos\frac{2\pi}{3}+i\sin\frac{2\pi}{3}\bigr)$.
\end{ejemplo}

\section{El producto en forma polar}

\begin{teorema}[Producto y cociente en forma polar]\label{teo:producto-polar}
Sean $z_1=r_1(\cos\theta_1+i\sin\theta_1)$ y $z_2=r_2(\cos\theta_2+i\sin\theta_2)$. Entonces
\[
z_1z_2=r_1r_2\bigl(\cos(\theta_1+\theta_2)+i\sin(\theta_1+\theta_2)\bigr),
\]
y, si $z_2\ne0$, $\dfrac{z_1}{z_2}=\dfrac{r_1}{r_2}\bigl(\cos(\theta_1-\theta_2)+i\sin(\theta_1-\theta_2)\bigr)$.
\end{teorema}

\begin{proof}
Desarrollando el producto y aplicando el \cref{cor:suma},
\[
(\cos\theta_1+i\sin\theta_1)(\cos\theta_2+i\sin\theta_2)=(\cos\theta_1\cos\theta_2-\sin\theta_1\sin\theta_2)+i(\sin\theta_1\cos\theta_2+\cos\theta_1\sin\theta_2),
\]
que es $\cos(\theta_1+\theta_2)+i\sin(\theta_1+\theta_2)$. Para el cociente, el número $w=\frac{r_1}{r_2}(\cos(\theta_1-\theta_2)+i\sin(\theta_1-\theta_2))$ cumple $wz_2=z_1$ por la primera parte, y el cociente es único porque $\C$ es un cuerpo.
\end{proof}

\begin{corolario}\label{cor:arg-producto}
Para $z,w\ne0$: $\abs{zw}=\abs{z}\abs{w}$ y $\arg(zw)=\arg z+\arg w$, en el sentido de que la suma de un argumento de $z$ y uno de $w$ es un argumento de $zw$. Para los argumentos principales, la igualdad solo vale salvo un múltiplo de $2\pi$.
\end{corolario}

\emph{Interpretación geométrica.} Fijado $w=\rho(\cos\alpha+i\sin\alpha)$, la aplicación $z\mapsto wz$ multiplica el módulo de cada punto por $\rho$ y suma $\alpha$ a su argumento: es la composición de un giro de ángulo $\alpha$ alrededor del origen y una homotecia de razón $\rho$. Esta lectura, que Needham sitúa en el centro de su exposición~\cite[cap.~1]{needham}, convierte en evidentes muchas identidades: por ejemplo, multiplicar por $i$ es girar un ángulo recto, y por eso $i^{2}=-1$ (dos giros rectos forman media vuelta).

\section{La exponencial imaginaria y la fórmula de Euler}

\begin{definicion}[Exponencial imaginaria]\label{def:euler}
Para $\theta\in\R$ se define
\[
e^{i\theta}=\cos\theta+i\sin\theta,
\]
y, para $z=x+iy\in\C$, $e^{z}=e^{x}e^{iy}=e^{x}(\cos y+i\sin y)$.
\end{definicion}

La notación no es arbitraria. Si se sustituye formalmente $x=i\theta$ en la serie $e^{x}=\sum_{n\ge0}x^{n}/n!$ y se agrupan los términos pares e impares usando $i^{2}=-1$, se obtienen exactamente las series de $\cos\theta$ y de $i\sin\theta$; las notas de Cambridge toman la serie como definición de $e^{z}$ y deducen la fórmula anterior~\cite[§1.2]{cambvm}. Las notas de 18.03 la justifican de otro modo: $f(\theta)=\cos\theta+i\sin\theta$ es la única función con $f'=if$ y $f(0)=1$, la misma ecuación que caracteriza a $e^{i\theta}$~\cite[§6.2]{mit1803}. Ambos argumentos requieren cálculo; aquí basta la definición y sus propiedades algebraicas.

\begin{proposicion}[Propiedades de $e^{i\theta}$]\label{prop:exp}
Para todos $\alpha,\beta,\theta\in\R$ y $n\in\Z$:
\begin{enumerate}
  \item $e^{i\alpha}e^{i\beta}=e^{i(\alpha+\beta)}$ y, más en general, $e^{z}e^{w}=e^{z+w}$ para $z,w\in\C$;
  \item $\abs{e^{i\theta}}=1$ y $\overline{e^{i\theta}}=e^{-i\theta}=1/e^{i\theta}$;
  \item $e^{i\theta}=1$ si y solo si $\theta\in2\pi\Z$; en particular, $e^{i\pi}=-1$ (identidad de Euler) y $e^{2\pi i}=1$;
  \item $\cos\theta=\dfrac{e^{i\theta}+e^{-i\theta}}{2}$ y $\sin\theta=\dfrac{e^{i\theta}-e^{-i\theta}}{2i}$.
\end{enumerate}
\end{proposicion}

\begin{proof}
(i) La primera igualdad es el \cref{teo:producto-polar} con $r_1=r_2=1$. Para $z=x+iy$ y $w=u+iv$, $e^{z}e^{w}=e^{x}e^{u}e^{iy}e^{iv}=e^{x+u}e^{i(y+v)}=e^{z+w}$, usando la propiedad de la exponencial real. (ii) $\abs{e^{i\theta}}^{2}=\cos^{2}\theta+\sin^{2}\theta=1$; $\overline{e^{i\theta}}=\cos\theta-i\sin\theta=\cos(-\theta)+i\sin(-\theta)=e^{-i\theta}$, y $e^{i\theta}e^{-i\theta}=e^{0}=1$ por (i). (iii) $e^{i\theta}=1$ equivale a $P(\theta)=P(0)$, y se aplica el \cref{teo:simetrias}(ii). (iv) Sumando y restando $e^{i\theta}=\cos\theta+i\sin\theta$ y $e^{-i\theta}=\cos\theta-i\sin\theta$.
\end{proof}

Con esta notación, todo $z\ne0$ se escribe en \emph{forma exponencial} $z=re^{i\theta}$, con $r=\abs{z}$ y $\theta$ un argumento, y el \cref{teo:producto-polar} se resume en $r_1e^{i\theta_1}\cdot r_2e^{i\theta_2}=r_1r_2e^{i(\theta_1+\theta_2)}$. El conjugado es $\overline{re^{i\theta}}=re^{-i\theta}$.

\section{El teorema de De Moivre}

\begin{teorema}[De Moivre]\label{teo:demoivre}
Para todo $\theta\in\R$ y todo $n\in\Z$,
\[
(\cos\theta+i\sin\theta)^{n}=\cos n\theta+i\sin n\theta,\qquad\text{es decir,}\qquad (e^{i\theta})^{n}=e^{in\theta}.
\]
\end{teorema}

\begin{proof}
Para $n\ge0$, por inducción. Si $n=0$, ambos miembros valen~1. Si se cumple para $n$, entonces $(e^{i\theta})^{n+1}=(e^{i\theta})^{n}e^{i\theta}=e^{in\theta}e^{i\theta}=e^{i(n+1)\theta}$ por la \cref{prop:exp}(i). Para $n=-m<0$, por la \cref{prop:exp}(ii), $(e^{i\theta})^{-m}=\bigl((e^{i\theta})^{m}\bigr)^{-1}=(e^{im\theta})^{-1}=e^{-im\theta}=e^{in\theta}$.
\end{proof}

\begin{corolario}
Si $z=re^{i\theta}$, entonces $z^{n}=r^{n}e^{in\theta}$ para todo $n\in\Z$ (para $n<0$, con $z\ne0$).
\end{corolario}

\begin{ejemplo}[Potencias y ángulo triple]\label{ej:demoivre}
(a) $1+i=\sqrt2\,e^{i\pi/4}$, luego $(1+i)^{8}=(\sqrt2)^{8}e^{2\pi i}=16$. (b) Desarrollando $(\cos\theta+i\sin\theta)^{3}$ por el binomio e igualando con $\cos3\theta+i\sin3\theta$,
\begin{align*}
\cos3\theta&=\cos^{3}\theta-3\cos\theta\sin^{2}\theta=4\cos^{3}\theta-3\cos\theta,\\
\sin3\theta&=3\cos^{2}\theta\sin\theta-\sin^{3}\theta=3\sin\theta-4\sin^{3}\theta,
\end{align*}
que son las fórmulas del \cref{eje:2-triple}, obtenidas ahora sin fórmulas de la suma.
\end{ejemplo}

\section{Raíces \texorpdfstring{$n$}{n}-ésimas y raíces de la unidad}

\begin{teorema}[Raíces $n$-ésimas]\label{teo:raices}
Sean $w=\rho e^{i\varphi}\ne0$ y $n\ge1$. La ecuación $z^{n}=w$ tiene exactamente $n$ soluciones distintas,
\[
z_k=\rho^{1/n}\,e^{i(\varphi+2\pi k)/n},\qquad k=0,1,\dots,n-1,
\]
donde $\rho^{1/n}$ es la raíz $n$-ésima real positiva de $\rho$.
\end{teorema}

\begin{proof}
Toda solución es no nula; se escribe $z=re^{i\theta}$. Por De Moivre, $z^{n}=r^{n}e^{in\theta}$, y $z^{n}=w$ equivale a $r^{n}=\rho$ (igualdad de módulos) y $e^{i(n\theta-\varphi)}=1$, es decir, por la \cref{prop:exp}(iii), $n\theta=\varphi+2\pi k$ con $k\in\Z$. Así $r=\rho^{1/n}$ y $\theta=(\varphi+2\pi k)/n$. Dos valores $k$ y $k'$ dan el mismo $z$ si y solo si $(\theta_k-\theta_{k'})=2\pi(k-k')/n\in2\pi\Z$, es decir, si $n$ divide a $k-k'$. Por tanto hay exactamente $n$ soluciones distintas, las correspondientes a $k=0,\dots,n-1$.
\end{proof}

Las $n$ raíces tienen el mismo módulo y argumentos separados $2\pi/n$: son los vértices de un polígono regular de $n$ lados centrado en el origen.

\begin{definicion}[Raíces de la unidad]
Las soluciones de $z^{n}=1$ son las \emph{raíces $n$-ésimas de la unidad}, $1,\omega,\omega^{2},\dots,\omega^{n-1}$ con $\omega=e^{2\pi i/n}$.
\end{definicion}

\begin{proposicion}\label{prop:suma-raices}
Si $n\ge2$, $\;1+\omega+\omega^{2}+\dots+\omega^{n-1}=0$.
\end{proposicion}

\begin{proof}
Como $\omega\ne1$ (porque $0<2\pi/n<2\pi$), la suma geométrica vale $\dfrac{\omega^{n}-1}{\omega-1}=\dfrac{1-1}{\omega-1}=0$.
\end{proof}

Geométricamente, la suma de los vectores que van del centro a los vértices de un polígono regular es nula: el conjunto es invariante por el giro de ángulo $2\pi/n$ (multiplicar la suma por $\omega$ la deja igual), y el único vector fijo por un giro no trivial es el cero.

\begin{figure}[ht]
\centering
\begin{tikzpicture}[scale=1.7,font=\small,>={Stealth[length=2mm]}]
  \draw[->] (-1.4,0) -- (1.5,0) node[right] {$\Real$};
  \draw[->] (0,-1.3) -- (0,1.4) node[above] {$\Imag$};
  \draw[gray] (0,0) circle (1);
  \draw[thick] (0:1) \foreach \k in {1,...,5} { -- (\k*60:1) } -- cycle;
  \foreach \k/\pos in {0/below right,1/above right,2/above left,3/below left,4/below left,5/below right}
    \fill (\k*60:1) circle (0.035) node[\pos] {$\omega^{\k}$};
  \draw[->,thick] (0,0) -- (60:1);
  \draw (0.3,0) arc (0:60:0.3);
  \node at (30:0.45) {$\frac{\pi}{3}$};
\end{tikzpicture}
\caption{Las raíces sextas de la unidad, $\omega^{k}=e^{ik\pi/3}$, son los vértices de un hexágono regular inscrito en la circunferencia unidad.}
\label{fig:raices}
\end{figure}

\section{Giros del plano como multiplicación por \texorpdfstring{$e^{i\alpha}$}{e\textasciicircum ia}}

\begin{proposicion}[Giros]\label{prop:giro-complejo}
Identifíquese $z=x+iy$ con el vector $(x,y)\T$. Para todo $\alpha\in\R$:
\begin{enumerate}
  \item la aplicación $z\mapsto e^{i\alpha}z$ corresponde a $(x,y)\T\mapsto R(\alpha)(x,y)\T$, donde $R(\alpha)$ es la matriz del \cref{cor:giros};
  \item conserva los módulos y los productos escalares: $\abs{e^{i\alpha}z}=\abs{z}$ y $\Real\bigl(e^{i\alpha}z\,\overline{e^{i\alpha}w}\bigr)=\Real(z\bar w)$;
  \item es biyectiva, con inversa $z\mapsto e^{-i\alpha}z$, y la composición de los giros de ángulos $\alpha$ y $\beta$ es el giro de ángulo $\alpha+\beta$.
\end{enumerate}
\end{proposicion}

\begin{proof}
(i) $e^{i\alpha}(x+iy)=(x\cos\alpha-y\sin\alpha)+i(x\sin\alpha+y\cos\alpha)$, cuyas coordenadas son las de $R(\alpha)(x,y)\T$. (ii) $\abs{e^{i\alpha}z}=\abs{e^{i\alpha}}\abs{z}=\abs{z}$, y $e^{i\alpha}z\,\overline{e^{i\alpha}w}=e^{i\alpha}e^{-i\alpha}z\bar w=z\bar w$ por la \cref{prop:exp}(ii). (iii) $e^{-i\alpha}(e^{i\alpha}z)=z$, y $e^{i\alpha}(e^{i\beta}z)=e^{i(\alpha+\beta)}z$ por la \cref{prop:exp}(i).
\end{proof}

La \cref{prop:giro-complejo}(iii) es la versión compleja del \cref{cor:giros}, y su demostración es más corta porque la regla de los exponentes ya contiene las fórmulas de la suma. Para girar alrededor de un punto $c$ distinto del origen se traslada, se gira y se deshace la traslación: $z\mapsto c+e^{i\alpha}(z-c)$.

\begin{corolario}[Productos escalares de vectores girados]\label{cor:rope}
Sean $z_1,z_2\in\C$, $\theta\in\R$ y $m,n\in\Z$. Entonces
\[
\Real\bigl(z_1e^{im\theta}\,\overline{z_2e^{in\theta}}\bigr)=\Real\bigl(z_1\bar z_2\,e^{i(m-n)\theta}\bigr),
\]
que solo depende de $m$ y $n$ a través de la diferencia $m-n$.
\end{corolario}

\begin{proof}
$\overline{z_2e^{in\theta}}=\bar z_2e^{-in\theta}$ por la \cref{prop:exp}(ii), y $e^{im\theta}e^{-in\theta}=e^{i(m-n)\theta}$ por la \cref{prop:exp}(i).
\end{proof}

\begin{observacion}[Errores frecuentes]
(a) Tomar $\arctan(b/a)$ como argumento sin mirar el cuadrante: para $-1-i$ daría $\pi/4$ en lugar de $-3\pi/4$. (b) Aplicar $\Arg(zw)=\Arg z+\Arg w$ a los argumentos principales: con $z=w=-1$ el primer miembro vale $0$ y el segundo $2\pi$. (c) Dar una sola raíz $n$-ésima: $z^{n}=w$ tiene siempre $n$ soluciones distintas si $w\ne0$. (d) Confundir $e^{i\theta}=e^{i\varphi}$ con $\theta=\varphi$: solo implica $\theta-\varphi\in2\pi\Z$.
\end{observacion}

\begin{observacion}[Conexión con el curso]
El \cref{cor:rope} es la proposición de la forma compleja de RoPE (módulo~14, tema~4): si cada par de coordenadas de la consulta del token $m$ se multiplica por $e^{im\theta}$ y el de la clave del token $n$ por $e^{in\theta}$, su producto escalar solo depende de la distancia $m-n$. Es la formulación del artículo original~\cite[§3.2]{roformer}. Por otra parte, la matriz $R(\alpha)$ no tiene autovalores reales si $\sin\alpha\ne0$, pero sí complejos, $e^{\pm i\alpha}$ (\cref{eje:6-autovalores}); el módulo~3 (tema~3) presenta el caso $\alpha=\pi/2$, con autovalores $\pm i$, como ejemplo de matriz real sin direcciones invariantes.
\end{observacion}

\begin{lecturas}
\item MIT 18.03, \emph{Supplementary Notes}~\cite{mit1803}, §6 «The complex exponential» (6.2 a 6.5).
\item Cambridge, \emph{Vectors and Matrices}~\cite{cambvm}, §1.2 «Complex exponential function», §1.3 «Roots of unity» y §1.5 «De Moivre's theorem».
\item Needham~\cite{needham}, cap.~1 «Geometry and Complex Arithmetic».
\item OpenStax, \emph{Algebra and Trigonometry}~\cite{osat}, §10.5 «Polar Form of Complex Numbers».
\end{lecturas}

\section*{Ejercicios del tema 6}
\addcontentsline{toc}{section}{Ejercicios del tema 6}

\begin{ejercicio}{1}{eje:6-exponenciales}
Escriba $1$, $i$, $-1$, $1+i$ y $\frac{1+\sqrt3\,i}{2}$ en forma exponencial $re^{i\theta}$.
\fuente{adaptado de MIT 18.03, \emph{Supplementary Notes}~\cite{mit1803}, ejercicio~6.3.1.}
\end{ejercicio}

\begin{ejercicio}{1}{eje:6-vf}
Justifique si son verdaderas o falsas: (a) $e^{i\pi}=-1$; (b) $\abs{e^{i\theta}}=1$ para todo $\theta$ real; (c) $\Arg(zw)=\Arg z+\Arg w$ para todos $z,w\ne0$; (d) si $e^{i\theta}=e^{i\varphi}$, entonces $\theta=\varphi$; (e) la conjugación $z\mapsto\bar z$ es un giro del plano.
\end{ejercicio}

\begin{ejercicio}{1}{eje:6-cuadrante}
Sin hacer cálculos largos, indique el cuadrante de $z=(1+i)^{5}$ y de $w=(1+i)/(1-i)^{2}$, razonando solo con los argumentos.
\end{ejercicio}

\begin{ejercicio}{2}{eje:6-polar}
Escriba en forma polar, con el argumento principal, los números $-1+i$, $-\sqrt3-i$, $-2i$ y $3-4i$ (para este último, dé el argumento con cuatro decimales).
\fuente{adaptado de OpenStax, \emph{Algebra and Trigonometry}~\cite{osat}, §10.5, ejercicios de la sección sobre la forma polar.}
\end{ejercicio}

\begin{ejercicio}{2}{eje:6-producto-polar}
Sean $z_1=10e^{i\pi/6}$ y $z_2=6e^{i\pi/3}$. Calcule $z_1z_2$ y $z_1/z_2$ en forma polar y binómica. Haga lo mismo con $z_1=12e^{i\,55^{\circ}}$ y $z_2=3e^{i\,18^{\circ}}$, dejando el resultado en forma polar.
\fuente{adaptado de OpenStax, \emph{Algebra and Trigonometry}~\cite{osat}, §10.5, ejercicios de la sección sobre productos y cocientes en forma polar.}
\end{ejercicio}

\begin{ejercicio}{2}{eje:6-potencias}
Calcule $\bigl(\frac{1+\sqrt3\,i}{2}\bigr)^{3}$ y $(1+i)^{4}$ en forma polar y compare con el cálculo en forma binómica.
\fuente{adaptado de MIT 18.03, \emph{Supplementary Notes}~\cite{mit1803}, ejercicios~5.1.1 y~6.4.1.}
\end{ejercicio}

\begin{ejercicio}{2}{eje:6-tripos}
Calcule $\dfrac{(1+i)^{2019}}{(1-i)^{2017}}$.
\fuente{adaptado de Cambridge, \emph{Mathematical Tripos} Part IA, Paper~1, 2019~\cite{tripos2019}, pregunta~1C(b).}
\end{ejercicio}

\begin{ejercicio}{2}{eje:6-raices}
(a) Escriba las raíces octavas de la unidad en forma binómica. (b) Calcule las raíces cúbicas de $-8$ y represéntelas. (c) Calcule las raíces sextas de $-i/64$ en forma exponencial.
\fuente{adaptado de MIT 18.03, \emph{Supplementary Notes}~\cite{mit1803}, ejercicios~6.5.1 y~6.5.2, y de Cambridge, \emph{Mathematical Tripos} Part IA, Paper~1, 2013~\cite{tripos2013}, pregunta~1C(a).}
\end{ejercicio}

\begin{ejercicio}{2}{eje:6-polinomios}
Encuentre todas las raíces complejas de $x^{3}+1$, $x^{4}-2x^{2}+1$ y $x^{4}+4$, y factorice $x^{4}+4$ como producto de dos polinomios reales de grado~2.
\fuente{adaptado de MIT 18.03, \emph{Supplementary Notes}~\cite{mit1803}, ejercicio~6.5.3 (se ha sustituido el segundo polinomio).}
\end{ejercicio}

\begin{ejercicio}{2}{eje:6-giros}
(a) Describa geométricamente la multiplicación por $i$, por $-1$, por $2i$ y por $\frac{1+i}{\sqrt2}$. (b) Gire el punto $(2,1)$ un ángulo $\pi/2$ y un ángulo $\pi/3$ alrededor del origen usando números complejos, y compruebe el primero con la matriz $R(\pi/2)$. (c) Gire $3+i$ un ángulo $\pi/2$ alrededor de $c=1+i$.
\end{ejercicio}

\begin{ejercicio}{2}{eje:6-rope}
Sean $z_1=1+2i$ (un par de coordenadas de una consulta) y $z_2=3+i$ (el par correspondiente de una clave), $\theta=\pi/6$, $m=3$ y $n=1$. Calcule $\Real\bigl(z_1e^{im\theta}\,\overline{z_2e^{in\theta}}\bigr)$ de dos maneras: girando primero cada número y usando el \cref{cor:rope}. Compruebe que se obtiene lo mismo con $m=5$ y $n=3$.
\end{ejercicio}

\begin{ejercicio}{3}{eje:6-suma-euler}
Suponga conocidas la fórmula de Euler y la propiedad $e^{z+w}=e^{z}e^{w}$ para todos $z,w\in\C$ (por ejemplo, demostrada a partir de la serie de potencias, como en~\cite[§1.2]{cambvm}). Deduzca de ellas las fórmulas del coseno y del seno de una suma, sin usar el \cref{cor:suma}.
\fuente{adaptado de MIT 18.03, \emph{Supplementary Notes}~\cite{mit1803}, ejercicio~6.4.2.}
\end{ejercicio}

\begin{ejercicio}{3}{eje:6-cos5}
(a) Usando el teorema de De Moivre y el binomio de Newton, demuestre que $\cos5\theta=16\cos^{5}\theta-20\cos^{3}\theta+5\cos\theta$. (b) Deduzca que $\cos^{2}\frac{\pi}{10}=\frac{5+\sqrt5}{8}$.
\fuente{adaptado de Cambridge, \emph{Vectors and Matrices}~\cite{cambvm}, §1.5, ejemplo del desarrollo de $\cos5\theta$.}
\end{ejercicio}

\begin{ejercicio}{3}{eje:6-potencias-raices}
Sean $n\ge2$ y $\omega=e^{2\pi i/n}$. Demuestre que, para todo $j\in\Z$,
\[
\sum_{k=0}^{n-1}\omega^{jk}=\begin{cases}n&\text{si $n$ divide a $j$},\\ 0&\text{en otro caso},\end{cases}
\]
y deduzca que $\sum_{k=0}^{n-1}\cos\frac{2\pi k}{n}=\sum_{k=0}^{n-1}\sin\frac{2\pi k}{n}=0$.
\fuente{adaptado de Cambridge, \emph{Vectors and Matrices}~\cite{cambvm}, §1.3, proposición sobre la suma de las raíces de la unidad.}
\end{ejercicio}

\begin{ejercicio}{3}{eje:6-suma-cosenos}
Sea $x\notin2\pi\Z$. Sumando la progresión geométrica $\sum_{k=0}^{n}e^{ikx}$ y tomando la parte real, demuestre que
\[
\sum_{k=0}^{n}\cos kx=\frac{\cos\frac{nx}{2}\,\sin\frac{(n+1)x}{2}}{\sin\frac{x}{2}} .
\]
Indicación: $e^{i\varphi}-1=e^{i\varphi/2}\cdot 2i\sin\frac{\varphi}{2}$.
\end{ejercicio}

\begin{ejercicio}{3}{eje:6-autovalores}
Sea $\alpha\in\R$ con $\sin\alpha\ne0$. Demuestre que la matriz $R(\alpha)$ no tiene autovalores reales, que sus autovalores complejos son $e^{i\alpha}$ y $e^{-i\alpha}$, y que $(1,-i)\T$ y $(1,i)\T$ son autovectores respectivos. Relacione el resultado con la \cref{prop:giro-complejo}(i).
\end{ejercicio}
